% AI-produced by GPT-6 Astra (OpenAI, ChatGPT Pro mode), 19 November 2026.
% Method/special-model result only. No movement credit on original A.5.
\subsection{Coherent approximation for one real-definable cell-like surface map}

We prove a finite height-leveling lemma, give the elementary dimension-three
link argument needed below, and apply them to an individual definable cell-like
map between closed surfaces over the real numbers. The first lemma works over
every real closed field. The last theorem uses classical topology over
$\mathbb R$ and is not asserted over arbitrary real closed fields or uniformly
in additional parameters. None of these claims closes the remaining A.5 gate.

\paragraph{Original-source inputs and precise locators.}
Rourke--Sanderson~\cite[Lemma 4.23, p.~58]{RourkeSanderson1972} levels a PL collar in a product by matching derived
subdivisions whose new vertices lie at a common positive height. We use that
finite construction for two possibly different nonnegative PL height functions
and supply the details. The same original gives the compact PL collar theorem
as Theorem 2.25, p.~24; the closed PL surface classification is used explicitly
in the remarks following Theorem 5.5, p.~63; the dimension-two PL Schoenflies
conclusion is recorded in the discussion of Problem 3.37, p.~47. Full original:
\url{https://webhomes.maths.ed.ac.uk/~v1ranick/papers/rsplx.pdf}.

Shiota~\cite[Theorem 2.2(1), p.~6]{A5-Shiota} gives exact definable
triangulation of a scalar function, with no reparametrization of its target.
Section 2 explicitly works over a real closed field. Full original:
\url{https://arxiv.org/pdf/1002.1508}.

Siebenmann~\cite[Approximation Theorem A, p.~271]{Siebenmann1972Approximation} applies to closed surfaces by condition (a).
Its Complement, p.~272, gives a continuous family of homeomorphisms with the
cell-like map as its endpoint. Full original:
\url{https://webhomes.maths.ed.ac.uk/~v1ranick/papers/siebappr.pdf}.
Both printed pp.~271--272 were visually reread from the original scan during
this audit. We use the continuous endpoint family only to establish that the
mapping cylinder is a topological three-manifold; the family is not claimed
to be definable.

\paragraph{Lemma 1 (finite leveling of a collared zero set).}
Let $R$ be a real closed field, let $Y$ be a compact polyhedron, and let
$g:Y\to[0,M]$ be PL. Set $L=g^{-1}(0)$ and suppose that $L$ admits a PL collar
in $Y$. Then there are $\epsilon>0$ and a PL homeomorphism
\[
 C:L\times[0,\epsilon]\longrightarrow g^{-1}([0,\epsilon])
 \quad\text{such that}\quad
 C(z,0)=z,\qquad g(C(z,s))=s.
 \tag{A5.16.1}
\]
In particular the given scalar height, rather than an unspecified collar
coordinate, is preserved.

\paragraph{Proof.}
Choose a PL collar $c:L\times[0,1]\to Y$, with $c(z,0)=z$, and write
$Y'=c(L\times[0,1])$. The image of $L\times[0,1)$ is an open neighborhood
of $L$ in $Y$. Its complement is a compact polyhedron disjoint from $L$,
so $g$ has a positive lower bound there. Thus a sufficiently small positive
$\epsilon$ has $g^{-1}([0,\epsilon])\subset Y'$.

Triangulate the finite PL diagram consisting of $c$, the domain height
$t:L\times[0,1]\to[0,1]$, and $g|Y'$. This gives finite complexes $K,J$
and a simplicial isomorphism
\[
 c:K\longrightarrow J,
 \qquad |K|=L\times[0,1],\quad |J|=Y',
\]
such that $t$ is linear on every simplex of $K$ and $g$ is linear on every
simplex of $J$. The zero vertices of $t$ correspond exactly to the zero
vertices of $g$, because both zero sets correspond to $L$ under $c$.
Shrink $\epsilon$ further so that $0<\epsilon<1$ and $\epsilon$ is less
than every positive number in the finite lists
\[
 \{t(v):v\in K^{(0)},\ t(v)>0\},\qquad
 \{g(c(v)):v\in K^{(0)},\ t(v)>0\}.
 \tag{A5.16.2}
\]

Call an original simplex of $K$ mixed if it has both a zero vertex and a
positive vertex. For each mixed simplex $\sigma$, its relative interior
meets $t^{-1}(\epsilon)$: positive barycentric coefficients may be chosen
whose weighted height is $\epsilon$, since the minimum vertex height is
zero and the maximum is greater than $\epsilon$. Similarly the relative
interior of $c(\sigma)$ meets $g^{-1}(\epsilon)$.

In decreasing order of dimension, stellarly subdivide every original mixed
simplex $\sigma$ at a point
\[
 a_\sigma\in\operatorname{relint}(\sigma)\cap t^{-1}(\epsilon),
\]
and subdivide $c(\sigma)$ at an independently chosen point
\[
 b_\sigma\in\operatorname{relint}(c(\sigma))\cap g^{-1}(\epsilon).
\]
Each original face is still present when it is to be starred: previous stars
were in relative interiors of larger faces or in other faces of the same
dimension. Paired stellar subdivisions therefore produce combinatorially
isomorphic complexes $K'$ and $J'$. Their canonical simplicial isomorphism
$c':K'\to J'$ sends each old vertex $v$ to $c(v)$ and each $a_\sigma$ to
$b_\sigma$, and fixes $L$ pointwise.

Every new vertex has height $\epsilon$. Every old positive vertex has height
greater than $\epsilon$ in the respective complex. No simplex of $K'$ can
contain both an old zero vertex and an old positive vertex. Indeed, these two
vertices would span an edge of $K'$; stellar subdivisions introduce no new
edges between old vertices, while every original zero--positive edge was
itself starred. The identical assertion holds in $J'$.

It follows that the sublevel set $t^{-1}([0,\epsilon])$ is exactly the
subcomplex of $K'$ with vertex heights $0$ or $\epsilon$. To see this
directly, a simplex containing a zero vertex contains no old positive vertex,
so all its heights are at most $\epsilon$. A simplex containing an old
positive vertex has no zero vertex, and its intersection with the sublevel
set is precisely the face spanned by its height-$\epsilon$ vertices.
The same description holds for $g^{-1}([0,\epsilon])$ in $J'$.

The isomorphism $c'$ identifies these two subcomplexes. On their vertices it
preserves height, since the only heights are $0$ and $\epsilon$; linearity
on each simplex then gives $g\circ c'=t$ throughout the band. Restricting
$c'$ gives (A5.16.1). All operations are finite PL operations, so the argument
works over $R$ and does not require a mesh whose size grows as the height
tends to zero. \hfill$\square$

\paragraph{Lemma 2 (the finite link argument in dimension three).}
Let $K$ be a finite simplicial complex over $\mathbb R$ whose underlying
space is a topological three-manifold, possibly with boundary. Assume that
the topological boundary is the underlying space of a subcomplex $B$.
Then $K$ is a PL three-manifold with boundary $B$.

\paragraph{Proof.}
First $K$ is pure of dimension three. An interior point of a maximal simplex
has a neighborhood of the dimension of that simplex; manifold dimension
and the existence of interior points in every manifold neighborhood force
every maximal simplex to have dimension three.

For a $k$-simplex $\sigma$ and a point $x\in\operatorname{relint}(\sigma)$,
the standard local homology formula is
\[
 H_i(|K|,|K|\setminus\{x\};\mathbb Z)
 \cong
 \widetilde H_{i-k-1}(\operatorname{lk}_K\sigma;\mathbb Z).
 \tag{A5.16.3}
\]
At an interior manifold point the left side is $\mathbb Z$ for $i=3$ and
zero otherwise. At a boundary point it is zero in every degree. Since $B$
is a subcomplex, each open simplex consists entirely of interior points or
entirely of boundary points.

If $\sigma$ is a two-simplex, its link is a nonempty finite set. Equation
(A5.16.3) says that the link has two points in the interior case and one point
in the boundary case. Thus every triangle has respectively two or one
incident tetrahedra.

If $\sigma$ is an edge, its link is a finite pure graph. Each vertex of this
graph corresponds to a triangle containing $\sigma$, so its degree is two
or one by the preceding paragraph. Equation (A5.16.3) makes the graph connected;
its first homology is $\mathbb Z$ in the interior case and zero in the
boundary case. A connected finite graph with all degrees one or two is
therefore a circle or an interval respectively. These are PL circles and
PL intervals, with their given finite triangulations.

For a vertex $v$, the link $S=\operatorname{lk}_K(v)$ is a compact pure
two-complex. The link of any vertex of $S$ is the link in $K$ of the
corresponding edge through $v$, hence is a circle or an interval. Therefore
$S$ is a compact PL surface, possibly with boundary. Equation (A5.16.3) makes
$S$ connected. If $v$ is interior, it has the integral homology of $S^2$.
The classification of compact PL surfaces makes it a PL two-sphere.

If $v$ is a boundary vertex, $S$ is acyclic. An acyclic compact connected
surface cannot be closed: an orientable closed surface has nonzero $H_2$,
and a nonorientable closed surface has nonzero integral $H_1$. Thus
$\partial S\ne\varnothing$. To include the boundary count explicitly, an
orientable compact connected surface of genus $r$ with $b>0$ boundary
components has
\[
 H_1(S;\mathbb Z)\cong\mathbb Z^{2r+b-1},
\]
while a nonorientable surface with $r\ge1$ crosscaps and $b>0$ boundary
components has $H_1(S;\mathbb Z)\cong\mathbb Z^{r+b-1}$.
Vanishing of $H_1$ forces the orientable case $r=0,b=1$; the nonorientable
case is impossible. Hence $S$ is a PL disk, and its boundary is connected.

For completeness this disk conclusion can be obtained using only the
closed-surface classification quoted above. Double $S$ along its $b$
boundary circles. Mayer--Vietoris gives $H_2(DS)\cong\mathbb Z^b$ and
$H_1(DS)\cong\mathbb Z^{b-1}$ because $S$ is acyclic. The double is a
connected closed PL surface, so $b=1$ and $DS$ is a PL sphere. The embedded
boundary circle is locally flat, and the two-dimensional PL Schoenflies
theorem makes either complementary closure, including $S$, a PL disk.

Thus every vertex link is the requisite PL sphere or PL ball. This is the
combinatorial link criterion for $K$ to be a PL three-manifold. The boundary
triangles are exactly those with one incident tetrahedron, so the PL boundary
is exactly the given topological boundary $B$. No assertion that an arbitrary
topological homeomorphism is PL is used. \hfill$\square$

\begin{saresult}{A5.016}{Proved}{Coherent definable approximation of one cell-like closed-surface map}
Fix an o-minimal expansion of $\mathbb R$. Let $F\subset\mathbb R^m$ and
$G\subset\mathbb R^n$ be compact definable closed topological surfaces,
and let $q:F\to G$ be a continuous definable surjection whose point inverses
are cell-like in the classical topological sense. Then there are
$\epsilon>0$ and a jointly definable continuous map
\[
 k:(0,\epsilon]\times G\longrightarrow F
\]
such that each $k_s$ is a homeomorphism and
\[
 \lim_{s\downarrow0}\sup_{y\in G}\|q(k_s(y))-y\|=0.
 \tag{A5.16.4}
\]
Equivalently, $q$ is the endpoint of a definable continuous family
$h_s:F\to G$ with $h_0=q$ and $h_s$ a homeomorphism for every $s>0$.
Consequently the associated double mapping-cylinder projection is definably
locally trivial at its central parameter.

\end{saresult}
\begin{proof}
\textbf{First establish the topological three-manifold.}
If $F=G=\varnothing$, the unique empty homeomorphism gives the conclusion,
with uniform error defined to be zero. If exactly one surface is empty,
the assumed surjection does not exist. We therefore assume both nonempty.
The assertion may be checked on components. Cell-like fibres are connected,
so components of $F$ map bijectively to components of $G$; both have finitely
many components. Assume the surfaces are connected.

By Siebenmann's Theorem A and its Complement there is a continuous map
$(t,x)\mapsto a_t(x)$ on $[0,1]\times F$ such that $a_1=q$ and
$a_t:F\to G$ is a homeomorphism for $0\le t<1$. Its continuity on a compact
domain gives uniform convergence $a_t\to q$ as $t\uparrow1$. For $s>0$ put
$b_s=a_{1-s}^{-1}$. The inverse family is jointly continuous for positive
$s$, as follows by taking inverses of the compact-domain homeomorphism
$(s,x)\mapsto(s,a_{1-s}(x))$ on any closed positive $s$-interval. Moreover
\[
 \sup_{y\in G}\|q(b_s(y))-y\|
 =\sup_{x\in F}\|q(x)-a_{1-s}(x)\|\longrightarrow0.
 \tag{A5.16.5}
\]

Consider the compact definable mapping cylinder
\[
 C_q=\{(s,sx,q(x)):0<s\le1,\ x\in F\}
       \cup\{(0,0,y):y\in G\},\qquad \tau(s,z,y)=s.
 \tag{A5.16.6}
\]
The map from $[0,1]\times G$ to $C_q$ given for $s>0$ by
\[
 (s,y)\longmapsto(s,sb_s(y),q(b_s(y)))
\]
and at $s=0$ by $(0,y)\mapsto(0,0,y)$ is continuous by boundedness of
$F$ and (A5.16.5), and is bijective. Compactness and the Hausdorff property
make it a homeomorphism. Thus $C_q$ is a topological three-manifold whose
boundary is the disjoint union of the levels $\tau=0$ and $\tau=1$.
This use of a classical family has established only this topological fact.

\medskip\noindent\textbf{Replace the height collar by a finite definable one.}
Apply Shiota's exact scalar triangulation theorem to
$\tau:C_q\to[0,1]$. There are a finite compact polyhedron $P$ and a
definable homeomorphism $\pi:P\to C_q$ such that $\ell=\tau\circ\pi$ is
PL. Subdivide so that $\ell$ is linear on each simplex. Since its values
lie in $[0,1]$, both levels $L_0=\ell^{-1}(0)$ and $L_1=\ell^{-1}(1)$
are subcomplexes. They triangulate precisely the topological boundary.
Lemma 2 makes $P$ a PL three-manifold, with boundary $L_0\sqcup L_1$.

The compact PL boundary component $L_0$ has a PL collar. Apply Lemma 1 to
$\ell$ to obtain a PL height-preserving collar
\[
 C:L_0\times[0,\epsilon]\longrightarrow\ell^{-1}([0,\epsilon]).
\]
Use the definable identification $j:G\to L_0$ determined by
$\pi(j(y))=(0,0,y)$. Then
\[
 \Phi(s,y)=\pi(C(j(y),s))
\]
is a definable height-preserving homeomorphism from $[0,\epsilon]\times G$
onto $\tau^{-1}([0,\epsilon])$, with the identity on the zero fibre.
For $s>0$ write
\[
 \Phi(s,y)=(s,s k_s(y),q(k_s(y))).
 \tag{A5.16.7}
\]
Division of the middle coordinate by $s$ makes $k$ jointly definable and
continuous for $s>0$; each $k_s$ is a homeomorphism because $\Phi$ identifies
the height-$s$ fibres. Uniform continuity of the last coordinate on the
compact domain $[0,\epsilon]\times G$ gives (A5.16.4).

Set $h_s=k_s^{-1}$ for $s>0$ and $h_0=q$. The inverse family is jointly
definable and continuous for positive $s$, and
\[
 \sup_{x\in F}\|h_s(x)-q(x)\|
 =\sup_{y\in G}\|y-q(k_s(y))\|\longrightarrow0.
\]
Hence the extension at $s=0$ is jointly continuous. Reflecting the positive
height collar using $|s|$ supplies the claimed double-cylinder chart.
\end{proof}

\paragraph{What a parameter-preserving version would additionally need.}
The finite proof of Lemma 1 does not automatically preserve an extra map
$p:Y\to A$, even if a PL collar $c$ already preserves $p$. In a mixed edge,
the domain point at height $\epsilon$ and the image point at height
$\epsilon$ can require different barycentric fractions. If $p$ varies
along that edge, those choices need not have the same $p$-coordinate.

There is a direct counterexample to the unrestricted parameter-preserving
claim. On $Y=[-1,1]\times[0,1]$ put $p(a,t)=a$, $u=|a|$, and
\[
 g(a,t)=
 \begin{cases}
 3t,&0\le t\le u/2,\\
 2u-t,&u/2\le t\le u,\\
 t,&u\le t\le1.
 \end{cases}
 \tag{A5.16.8}
\]
The function is continuous and PL, positive exactly when $t>0$. The zero
set has the standard product PL collar over $a$. Each fixed $a$ admits
a height-preserving collar near its zero point. Nevertheless no uniform
collar near $(a,g)=(0,0)$ can preserve both $p$ and $g$: for every sufficiently
small $a>0$ and every $a<s<3a/2$, the fibre over $(a,s)$ has the three
distinct points
\[
 t=s/3,\qquad t=2a-s,\qquad t=s,
\]
whereas the central fibres have one point.

This counterexample is not a counterexample to the family mapping-cylinder
problem. In that problem every positive-height fibre is already the same
surface $F$; (A5.16.8) violates that condition. It shows precisely why a
parameter-preserving collar of the zero family, or a fibrewise PL manifold
structure alone, cannot justify parameter-preserving height leveling.

For a family $q_a:F\to G$, the needed conclusion remains a jointly definable
family $k(a,s,y)$ with $q_a k(a,s,y)\to y$ jointly at $s=0$. To use a finite
PL argument, one would need a parameter-compatible triangulation and a
leveling construction whose stellar choices preserve the parameter map, or
another genuine family theorem exploiting the special positive-height
product. The present argument does not supply either assertion.

\paragraph{Status.}
Lemma 1 closes the previously isolated prescribed-height collar substep.
Lemma 2 avoids an unjustified appeal to uniqueness of PL structures applied
to an arbitrary topological homeomorphism. A5.016 treats one cell-like
surface map over $\mathbb R$. It is a special-model result, not a reduction
or closure of the whole A.5 problem; no original-gate movement credit is
claimed. The uniform family and arbitrary-real-closed-field conclusions
remain unproved here.

\paragraph{Finite exact certificate.}
The accompanying checker
\path{check_surface_collar_parameter_fold.py}
verifies all three affine pieces in (A5.16.8), their two boundary matching identities, and the three distinct preimages with their region inequalities. It records 21 exact checks passing in 0.06886558799305931 seconds. These checks certify the displayed parameter-fold algebra only; the finite PL collar, link argument, and cell-like surface theorem are proved above.
